\section{Theorem 3.1 for projective space in characteristic zero}\label{sec:char0-blueprint}

This chapter proves the envelope property of $\PP^n$ (\cite[Theorem~3.1]{FHX}, Theorem~\ref{thm:pn-ep})
when $k$ is a countable algebraically closed field of characteristic $0$, the standing assumption of
\cite[\S5]{FHX}. For this case \cite{FHX} cites \cite[Theorem~5.20]{BJ22}, whose proof resolves the
singularities of a test configuration and uses the asymptotic multiplier ideals of its base ideals,
their subadditivity and their uniform global generation (Nadel vanishing). The route below is
different and uses none of these; Remark~\ref{rmk:c0-source} compares the two.
\begin{itemize}
\item As in characteristic $p$ (Remark~\ref{rmk:bj-pn-ep-assembly}), it suffices to show that
$\mathrm P_d(\phi)=\sup_m m^{-1}\phi_{J_m(\phi)}$ is $d$-psh for a PL obstacle $\phi\le0$. This follows from the
uniform estimate $m\,\mathrm P_d(\phi)\le\phi_{J_{m+2}(\phi)}+N_0$ (Proposition~\ref{prop:c0-key-inequality}),
which gives uniform convergence of $m^{-1}\phi_{J_m}$ (Theorem~\ref{thm:c0-envelope-pl}).
\item The field $k$ is embedded in $\mathbf C$, and the divisorial valuation rings over $k$ are extended to
$\mathbf C$ without changing orders (Lemma~\ref{lem:c0-base-change-dvr}).
\item The analytic input is Demailly's lower bound for Bergman kernels, obtained from the
Ohsawa--Takegoshi extension theorem at a point and globalized on $\mathbf C^{n+2}$ by H\"ormander's
$L^2$ estimate with polynomial weights (Proposition~\ref{prop:c0-l2-extension}). It produces, for
$J_q=J_q(\phi)$ and $c=m/q$, finitely many polynomials of $x$-degree $\le md+2$ which are square
integrable against $|J_q|^{-2c}$ and whose modulus dominates $|J_q|^{c}$
(Corollary~\ref{cor:c0-bergman-lower}).
\item The valuative input is the bound $\operatorname{ord}_V(f)>c\operatorname{ord}_V(g)-A(V)$ for polynomials $f$ which are
square integrable against $|g|^{-2c}$, at every divisorial valuation ring $V$
(Proposition~\ref{prop:c0-valuative-bound}). It is proved in coordinates at a general point of the
divisor of $V$; no resolution of singularities is used.
\end{itemize}

\begin{convention}\label{conv:c0}
\uses{conv:bj-pn, def:pn-div, def:pn-gauss}
In this chapter $k$ is a countable algebraically closed field of characteristic $0$, $n\ge1$, $N=n+2$, and
$z=(x_0,\dots,x_n,t)$, so that $R[t]=k[x,t]=k[z]$ is the coordinate ring of the affine cone over
$\PP^n\times\A^1$ and the discrete valuation rings of Definition~\ref{def:pn-div} are subrings of
$k(z)$. On $\mathbf C^N=\mathbf C^{n+1}_x\times\mathbf C_t$ we use the Lebesgue measure $\lambda$ and
$|z|^2=|x|^2+|t|^2$; $B(y,r)$ is the open ball. For a finite family $g=(g_1,\dots,g_s)$ of holomorphic
functions, not all identically zero, $|g|=(\sum_j|g_j|^2)^{1/2}$, $Z(g)=\{|g|=0\}$, and $\log|g|$ takes the
value $-\infty$ on $Z(g)$. For $c>0$ and $f$ holomorphic we say that \emph{$f$ is square integrable
against $|g|^{-2c}$} if $\int_K|f|^2|g|^{-2c}\,d\lambda<\infty$ for every compact $K$ (the integrand is
read as $0$ where $f=0$). A polynomial in $\mathbf C[x,t]$ is \emph{$x$-homogeneous} of degree $e$ if it
is homogeneous of degree $e$ in $x$, and \emph{bihomogeneous} if moreover it is a monomial in $t$; its
$x$-degree and $t$-degree are its degrees in $x$ and in $t$.
\end{convention}

\subsection{Divisorial valuation rings and extension of scalars}

\begin{definition}[Divisorial valuation rings]\label{def:c0-dvr}
\lean{EnvelopeProperty.BirationalDVR, EnvelopeProperty.BirationalDVR.ordIdeal}
\leanok
\uses{conv:c0}
\source{fhx-envelope-20261028:subsec:prelim-berkovich}
Let $F$ be a field and $N\ge1$. A \emph{divisorial valuation ring over $F$} (in the variables
$z_1,\dots,z_N$) is a discrete valuation ring $V$ with $F[z]\subseteq V\subseteq F(z)$ which is essentially of
finite type over $F$. We write $\operatorname{ord}_V$ for its normalized valuation, $\mathfrak m_V$ and $\kappa_V$ for its maximal
ideal and residue field, $\operatorname{ord}_V(g)=\min_j\operatorname{ord}_V(g_j)$ for a finite family of elements of $V$, and
$\operatorname{ord}_V(I)=\min_{f\in I}\operatorname{ord}_V(f)$ for an ideal $I$ of $F[z]$. A point $v\in(\PP^n)^\an$ is divisorial
(Definition~\ref{def:pn-div}) if and only if $\log|g|_{\sigma(v)}=-c\operatorname{ord}_V(g)$ on $k[z]$ for some
$c\in\Q_{>0}$ and some divisorial valuation ring $V$ over $k$ in the $N=n+2$ variables $z$.
\end{definition}

\begin{lemma}[Residue fields of divisorial valuation rings]\label{lem:c0-residue-trdeg}
\lean{EnvelopeProperty.Char0.BirationalDVR.trdeg_residueField, EnvelopeProperty.Char0.trdeg_residueField_add_one, EnvelopeProperty.Char0.height_eq_of_isMaximal_mvPolynomial}
\leanok
\uses{def:c0-dvr}
Let $V$ be a divisorial valuation ring over a field $F$ in $N$ variables. Then $\kappa_V$ is a finitely
generated field extension of $F$ of transcendence degree $N-1$.
\end{lemma}

\begin{proof}
Since $V$ is essentially of finite type, $V=E_P$ for a finitely generated $F$-subalgebra $E\subseteq V$
containing $z_1,\dots,z_N$ and $P=\mathfrak m_V\cap E$; then $K\coloneqq\operatorname{Frac}E=F(z)$ and $\kappa_V=\operatorname{Frac}(E/P)$ is
finitely generated over $F$. Let $\bar u_1,\dots,\bar u_r$ be a transcendence basis of $\kappa_V$ over $F$ and,
enlarging $E$, let $u_i\in E$ be lifts. A nonzero $h\in F[u]$ has nonzero residue $h(\bar u)$, so $h$ is a unit
of $V$: the $u_i$ are algebraically independent and $F'\coloneqq F(u)\subseteq V$. A uniformizer $\pi$ of $V$ is
transcendental over $F'$: an equation $\pi^e+a_1\pi^{e-1}+\dots+a_e=0$ with $a_i\in F'\subseteq V$ and $a_e\ne0$, a unit,
would give $a_e\in\pi V$. Hence $s\coloneqq\operatorname{trdeg}_{F'}K=N-r\ge1$.

Let $S=F[u]\setminus\{0\}$, $E'=S^{-1}E\subseteq V$ and $P'=S^{-1}P$. Then $E'$ is a finitely generated $F'$-domain with
fraction field $K$ and $E'_{P'}=V$, so $\operatorname{ht}P'=1$. The domain $E'/P'=S^{-1}(E/P)\subseteq\kappa_V$ is finitely
generated over $F'$ and algebraic over it ($\kappa_V$ is algebraic over $F(\bar u)$), hence finite dimensional, hence a
field: $P'$ is maximal. By Noether normalization there are $y_1,\dots,y_s\in E'$, algebraically independent over $F'$,
with $E'$ finite over $A_0=F'[y_1,\dots,y_s]$. The prime $\mathfrak m_0=P'\cap A_0$ is maximal, as $E'/P'$ is a field integral
over $A_0/\mathfrak m_0$.

Every maximal ideal $\mathfrak M$ of $F'[y_1,\dots,y_s]$ has height $\ge s$. By induction on $s$, write
$F'[y]=R_1[y_s]$ with $R_1=F'[y_1,\dots,y_{s-1}]$. The quotient $F'[y]/\mathfrak M$ is finite over $F'$ (Zariski's lemma), so
the domain $R_1/\mathfrak p$, $\mathfrak p=\mathfrak M\cap R_1$, is finite over $F'$ and $\mathfrak p$ is maximal. A chain
$\mathfrak q_0\subsetneq\dots\subsetneq\mathfrak q_{s-1}=\mathfrak p$ of primes of $R_1$ (induction) gives the chain
$\mathfrak q_0[y_s]\subsetneq\dots\subsetneq\mathfrak p[y_s]\subsetneq\mathfrak M$: the ideals $\mathfrak q[y_s]$ are prime since
$F'[y]/\mathfrak q[y_s]=(R_1/\mathfrak q)[y_s]$, the inclusions are strict on contracting to $R_1$, and $\mathfrak p[y_s]\ne\mathfrak M$ because
$(R_1/\mathfrak p)[y_s]$ is not a field.

Since $A_0$ is integrally closed and the domain $E'$ is integral over it, going down holds for $A_0\subseteq E'$. A chain
of primes of length $\ge s$ ending at $\mathfrak m_0$ therefore lifts to a chain of the same length ending at $P'$. Hence
$1=\operatorname{ht}P'\ge s\ge1$, so $s=1$ and $r=N-1$.
\end{proof}

\begin{lemma}[Tensor products over an algebraically closed field]\label{lem:c0-tensor-domain}
\lean{EnvelopeProperty.isDomain_tensorProduct_of_isAlgClosed, EnvelopeProperty.eq_zero_of_forall_algHom_eq_zero}
\leanok
Let $k$ be an algebraically closed field and $A$, $B$ $k$-algebras which are domains. Then $A\otimes_kB$ is a domain.
\end{lemma}

\begin{proof}
$A\otimes_kB\ne0$. Let $\alpha,\beta\in A\otimes_kB$ with $\alpha\beta=0$. Their $A$-coefficients lie in a finitely generated
subalgebra $A_0\subseteq A$, and $A_0\otimes_kB\to A\otimes_kB$ is injective ($k$ is a field), so we may assume $A$ finitely
generated. Write $\alpha=\sum_ia_i\otimes b_i$ and $\beta=\sum_jc_j\otimes d_j$ with $(b_i)$ and $(d_j)$ each linearly independent
over $k$. For a maximal ideal $\mathfrak m$ of $A$ we have $A/\mathfrak m=k$ (Nullstellensatz), and the $k$-algebra map
$A\otimes_kB\to B$, $a\otimes b\mapsto\bar ab$, sends $\alpha\beta$ to $(\sum_i\bar a_ib_i)(\sum_j\bar c_jd_j)=0$; as $B$ is a domain
and the $b_i$, $d_j$ are independent, all $\bar a_i$ or all $\bar c_j$ vanish. So every maximal ideal contains
$I_1=(a_i)_i$ or $I_2=(c_j)_j$, hence contains $I_1I_2$. The Jacobson radical of the finitely generated reduced
$k$-algebra $A$ is $0$, so $I_1I_2=0$, and $I_1=0$ or $I_2=0$ since $A$ is a domain; that is, $\alpha=0$ or $\beta=0$.
\end{proof}

\begin{lemma}[Extension of scalars of divisorial valuation rings]\label{lem:c0-base-change-dvr}
\lean{EnvelopeProperty.BirationalDVR.baseChange, EnvelopeProperty.BirationalDVR.irreducible_toBaseChange, EnvelopeProperty.BirationalDVR.addVal_toBaseChange_sum, EnvelopeProperty.BirationalDVR.ord_baseChange_sum, EnvelopeProperty.BirationalDVR.coneBaseChange, EnvelopeProperty.BirationalDVR.ord_coneBaseChange_sum}
\leanok
\uses{def:c0-dvr, lem:c0-tensor-domain}
\source{bj22-trivval:sec:Gauss}
Let $k$ be algebraically closed, $K'\supseteq k$ a field, and $V$ a divisorial valuation ring over $k$ in $N$ variables.
Then $V\otimes_kK'$ is a domain, $\mathfrak m_V\otimes_kK'$ is a prime ideal, and
$V_{K'}\coloneqq(V\otimes_kK')_{\mathfrak m_V\otimes K'}$ is a divisorial valuation ring over $K'$ in the same variables, with
uniformizer $\pi\otimes1$ for a uniformizer $\pi$ of $V$, such that
\[
\operatorname{ord}_{V_{K'}}\Bigl(\sum_\beta f_\beta\otimes b_\beta\Bigr)=\min_\beta\operatorname{ord}_V(f_\beta)
\]
for all finite families $f_\beta\in V$ and $b_\beta\in K'$ with the $b_\beta$ linearly independent over $k$.
\end{lemma}

\begin{proof}
$V\otimes_kK'$ and $\kappa_V\otimes_kK'=(V\otimes_kK')/(\mathfrak m_V\otimes K')$ are domains (Lemma~\ref{lem:c0-tensor-domain}), so the
ideal is prime. Write $V=S^{-1}E$ with $E$ a finitely generated $k$-algebra; then $V\otimes_kK'=S^{-1}(E\otimes_kK')$ is a
localization of a finitely generated $K'$-algebra, hence Noetherian, and $V_{K'}$ is a Noetherian local domain, essentially
of finite type over $K'$. Its maximal ideal is generated by $\pi\otimes1$, since $\mathfrak m_V=\pi V$, and $\pi\otimes1\ne0$ by
flatness of $K'$ over $k$; a Noetherian local domain whose maximal ideal is principal and nonzero is a discrete
valuation ring. As $E\subseteq k(z)$ and $k(z)\otimes_kK'$ is the localization of $K'[z]=k[z]\otimes_kK'$ at the nonzero elements
of $k[z]$, we get $K'[z]\subseteq V_{K'}\subseteq K'(z)$.
For the formula, write $f_\beta=\pi^{n_\beta}\varepsilon_\beta$ with $\varepsilon_\beta\in V^\times$ (the terms with $f_\beta=0$ may be dropped) and let
$n=\min n_\beta$. Then $\sum f_\beta\otimes b_\beta=(\pi\otimes1)^n\bigl(\sum_{n_\beta=n}\varepsilon_\beta\otimes b_\beta+(\pi\otimes1)\gamma\bigr)$ with
$\gamma\in V\otimes K'$. The elements $1\otimes b_\beta$ are linearly independent over $\kappa_V$ in $\kappa_V\otimes_kK'$, and the
$\bar\varepsilon_\beta$ are nonzero, so the bracket has nonzero image $\sum_{n_\beta=n}\bar\varepsilon_\beta\otimes b_\beta$ in $\kappa_V\otimes K'$, that is, it
is a unit of $V_{K'}$. Hence the order is $n$.
\end{proof}

\begin{corollary}[Orders and filtrations under extension of scalars]\label{cor:c0-base-change-orders}
\lean{EnvelopeProperty.BirationalDVR.mem_baseChange_ordIdeal_iff, EnvelopeProperty.BirationalDVR.ord_baseChange_one_tmul, EnvelopeProperty.BirationalDVR.iInf_ord_baseChange_map, EnvelopeProperty.BirationalDVR.mem_span_coneMap_ordIdeal_iff, EnvelopeProperty.BirationalDVR.ord_coneBaseChange_map, EnvelopeProperty.BirationalDVR.iInf_ord_coneBaseChange_map}
\leanok
\uses{lem:c0-base-change-dvr}
In the setting of Lemma~\ref{lem:c0-base-change-dvr}:
\begin{enumerate}
\item for a $k$-subspace $W\subseteq k[z]$ and $\mu\in\R$,
$\{G\in W\otimes_kK':\operatorname{ord}_{V_{K'}}(G)\ge\mu\}=\{f\in W:\operatorname{ord}_V(f)\ge\mu\}\otimes_kK'$, and the same holds for finite
intersections of such conditions attached to finitely many divisorial valuation rings $V_1,\dots,V_r$ over $k$;
\item for an ideal $I\subseteq k[z]$, $\operatorname{ord}_{V_{K'}}(IK'[z])=\operatorname{ord}_V(I)$; for $f\in k[z]$, $\operatorname{ord}_{V_{K'}}(f)=\operatorname{ord}_V(f)$.
\end{enumerate}
\end{corollary}

\begin{proof}
(1) Fix a $k$-basis $(b_\beta)$ of $K'$. Every $G\in W\otimes K'$ is uniquely $\sum_\beta f_\beta\otimes b_\beta$ with $f_\beta\in W$, and by
Lemma~\ref{lem:c0-base-change-dvr} $\operatorname{ord}(G)\ge\mu$ if and only if $\operatorname{ord}_V(f_\beta)\ge\mu$ for all $\beta$. For finitely many
$V_i$ the condition on $G$ is again that every coordinate $f_\beta$ satisfies all conditions. (2) If $I=(g_1,\dots,g_s)$, an
element $\sum h_jg_j$ of $IK'[z]$ has order $\ge\min_j\operatorname{ord}(g_j)$, and $\operatorname{ord}_{V_{K'}}(g_j)=\operatorname{ord}_V(g_j)$ by the formula with a single
term $b=1$.
\end{proof}

\begin{lemma}[Embedding in the complex numbers]\label{lem:c0-embedding}
\lean{EnvelopeProperty.nonempty_ringHom_complex}
\leanok
A countable algebraically closed field $k$ of characteristic $0$ admits a field embedding $\iota\colon k\to\mathbf C$.
\end{lemma}

\begin{proof}
Let $T$ be a transcendence basis of $k$ over $\Q$; it is countable. A transcendence basis $T'$ of $\mathbf C$ over $\Q$ is
uncountable, since the algebraic closure of $\Q(T'')$ is countable for countable $T''$. Choose an injection $T\to T'$;
it induces a field embedding $\Q(T)\to\Q(T')\subseteq\mathbf C$, which extends to the algebraic extension $k$ of $\Q(T)$ because
$\mathbf C$ is algebraically closed.
\end{proof}

\subsection{Coordinates at the generic point of a divisor}

In this subsection $F$ is a field of characteristic $0$ and $V$ a divisorial valuation ring over $F$ in $N$ variables, with
fraction field $K=F(z)$.

\begin{lemma}[Divisorial valuation rings as local rings of curves]\label{lem:c0-dvr-curve}
\lean{EnvelopeProperty.Char0.curveDVR, EnvelopeProperty.Char0.isDiscreteValuationRing_curveDVR, EnvelopeProperty.Char0.mem_curveDVR_iff, EnvelopeProperty.Char0.curveB, EnvelopeProperty.Char0.isDedekindDomain_curveB, EnvelopeProperty.Char0.exists_mul_mem_curveB, EnvelopeProperty.Char0.mem_of_mul_mem_curveB, EnvelopeProperty.Char0.curveBToResidue_surjective, EnvelopeProperty.Char0.isSeparable_curveField}
\leanok
\uses{lem:c0-residue-trdeg}
Let $u_2,\dots,u_N\in V$ lift a transcendence basis of $\kappa_V$ over $F$ (Lemma~\ref{lem:c0-residue-trdeg}) and let $\pi$ be a
uniformizer. Then $L\coloneqq F(u_2,\dots,u_N)\subseteq V$, $\pi$ is transcendental over $L$, $K$ is a finite separable extension of
$L(\pi)$, and $A\coloneqq L[\pi]_{(\pi)}$ is a discrete valuation ring with $A\subseteq V$ and $\pi A=\mathfrak m_V\cap A$. The integral closure
$B$ of $A$ in $K$ is a free $A$-module of rank $[K:L(\pi)]$ and a Dedekind domain contained in $V$, and $V=B_{\mathfrak n}$ for
$\mathfrak n=\mathfrak m_V\cap B$. The residue field $\kappa_V=B/\mathfrak n$ is finite over $L$, and $\mathfrak nB_{\mathfrak n}=\pi B_{\mathfrak n}$.
\end{lemma}

\begin{proof}
As in the proof of Lemma~\ref{lem:c0-residue-trdeg}, nonzero elements of $F[u]$ are units of $V$, so $L\subseteq V$, and $\pi$ is
transcendental over $L$. Since $\operatorname{trdeg}_FK=N$, $K$ is algebraic over $L(\pi)$, hence finite (it is finitely generated), and
separable (characteristic $0$). An element $Q(\pi)/S(\pi)$ of $A$ with $S(0)\ne0$ has $S(\pi)\equiv S(0)\in L^\times$ modulo
$\mathfrak m_V$, so $A\subseteq V$ and $\mathfrak m_V\cap A=\pi A$. As $A$ is a Noetherian integrally closed domain and $K/L(\pi)$ is separable,
$B$ is finite over $A$, torsion free, hence free over the principal ideal domain $A$, of rank $[K:L(\pi)]$; it is the
integral closure of a Dedekind domain in a finite separable extension, hence Dedekind. $B\subseteq V$ because $V$ is integrally
closed and contains $A$. The localization $B_{\mathfrak n}$ is a discrete valuation ring with fraction field $K$ and $V$ dominates
it; if $x\in V\setminus B_{\mathfrak n}$ then $x^{-1}\in\mathfrak nB_{\mathfrak n}\subseteq\mathfrak m_V$, contradicting $x\in V$. So $V=B_{\mathfrak n}$, and
$\mathfrak nB_{\mathfrak n}=\mathfrak m_V=\pi V$. Finally $\kappa_V$ is algebraic over $L$ (it has transcendence degree $N-1$ over $F$ and
contains the image of $L$) and finitely generated, hence finite.
\end{proof}

\begin{lemma}[A primitive coordinate]\label{lem:c0-dvr-primitive}
\lean{EnvelopeProperty.Char0.IsPrimitiveCoordinate, EnvelopeProperty.Char0.exists_isPrimitiveCoordinate}
\leanok
\uses{lem:c0-dvr-curve}
In the setting of Lemma~\ref{lem:c0-dvr-curve} there is $\theta\in B$ such that
\begin{enumerate}
\item $K=L(\pi)(\theta)$, and the residue $\xi=\bar\theta\in\kappa_V$ is nonzero and generates $\kappa_V$ over $L$;
\item the minimal polynomial $F_\theta$ of $\theta$ over $L(\pi)$ lies in $A[T]$, and its reduction $\bar F_\theta\in L[T]$ modulo $\pi$
has $\xi$ as a simple root: $\bar F_\theta=G\cdot T^{e}$ with $G$ the minimal polynomial of $\xi$ over $L$, separable, $G(0)\ne0$;
\item $V=A[\theta]_{\mathfrak n'}$ with $\mathfrak n'=\mathfrak m_V\cap A[\theta]$.
\end{enumerate}
\end{lemma}

\begin{proof}
Let $\mathfrak n=\mathfrak n_1,\mathfrak n_2,\dots,\mathfrak n_g$ be the maximal ideals of $B$ containing $\pi$, so that $\pi B=\prod_j\mathfrak n_j^{e_j}$ with $e_1=1$
(Lemma~\ref{lem:c0-dvr-curve}). The extension $\kappa_V/L$ is finite separable; choose a primitive element $\xi\ne0$ (replace a
primitive element $\xi_0$ by $\xi_0+1$ if $\xi_0=0$, by $1$ if $\kappa_V=L$). By the Chinese remainder theorem there is $\theta_1\in B$ with
$\theta_1\equiv\xi$ modulo $\mathfrak n_1$ and $\theta_1\in\mathfrak n_j$ for $j\ge2$. Let $I=\mathfrak n_1\cdots\mathfrak n_g\supseteq\pi B$, so $I$ spans $K$ over $L(\pi)$,
and let $W_1,\dots,W_h$ be the proper intermediate fields of $K/L(\pi)$ (finitely many, $K/L(\pi)$ being finite separable).
Some $a\in I$ lies in no $W_i$: otherwise let $W_1,\dots,W_{h'}$ be a minimal subfamily covering $I$; $h'\ge2$ since $I$ spans
$K$; choose $a_1\in I\setminus(W_2\cup\dots\cup W_{h'})$ (so $a_1\in W_1$) and $b\in I\setminus W_1$; the elements $b+\lambda a_1\in I$, $\lambda\in F$, are not
in $W_1$, so two of them with $\lambda\ne\lambda'$ lie in the same $W_i$, $2\le i\le h'$ ($F$ is infinite), whence $a_1\in W_i$, a contradiction. In the same way some $\theta=\theta_1+\lambda a$, $\lambda\in F$, lies in no $W_i$: two values in one $W_i$ would put $a$
in $W_i$. This $\theta$ satisfies (1), since $\theta\equiv\theta_1$ modulo $I$.

(2) $\theta$ is integral over the integrally closed domain $A$, so $F_\theta\in A[T]$, monic of degree $[K:L(\pi)]=\operatorname{rank}_AB$. As $\theta$
generates $K$, $F_\theta$ is the characteristic polynomial of multiplication by $\theta$ on $K$ over $L(\pi)$, hence on the free $A$-module
$B$. Reducing modulo $\pi$, $\bar F_\theta$ is the characteristic polynomial of $\theta$ on $B/\pi B\cong\prod_jB/\mathfrak n_j^{e_j}$ over $L$. On
$B/\mathfrak n_1=\kappa_V$ it is the minimal polynomial $G$ of the primitive element $\xi$; on $B/\mathfrak n_j^{e_j}$, $j\ge2$, $\theta$ is nilpotent and the
characteristic polynomial is a power of $T$. So $\bar F_\theta=G\,T^{e}$, $G$ is separable, and $G(0)\ne0$ because $G$ is irreducible and
$\xi\ne0$; thus $\xi$ is a simple root.

(3) $C\coloneqq A[\theta]\cong A[T]/(F_\theta)$, and $C/\pi C\cong L[T]/(\bar F_\theta)$. Lift $G$ and $T^e$ to monic $\widetilde G,H\in A[T]$; then
$F_\theta=\widetilde GH+\pi Q$ with $Q\in A[T]$. The prime $\mathfrak n'$ contains $\pi$ and corresponds to the factor $G$, so
$\mathfrak n'=(\pi,\widetilde G(\theta))$; $H(\theta)\notin\mathfrak n'$ since $\xi^e\ne0$. In $C_{\mathfrak n'}$, $0=F_\theta(\theta)$ gives
$\widetilde G(\theta)=-\pi Q(\theta)H(\theta)^{-1}$, so the maximal ideal is $\pi C_{\mathfrak n'}$. Hence $C_{\mathfrak n'}$ is a Noetherian local domain with
nonzero principal maximal ideal, a discrete valuation ring, with fraction field $L(\pi)(\theta)=K$, and $V$ dominates it
($C\subseteq B\subseteq V$, $\mathfrak n'=\mathfrak m_V\cap C$). As in Lemma~\ref{lem:c0-dvr-curve}, $V=C_{\mathfrak n'}$.
\end{proof}

\begin{lemma}[Partial derivatives and the log discrepancy]\label{lem:c0-dvr-jacobian}
\lean{EnvelopeProperty.Char0.partialDeriv, EnvelopeProperty.Char0.partialDeriv_mem, EnvelopeProperty.Char0.jacobian, EnvelopeProperty.Char0.jacobian_ne_zero, EnvelopeProperty.Char0.BirationalDVR.jacobian_mem_and_ne_zero, EnvelopeProperty.Char0.logDiscrepancy, EnvelopeProperty.Char0.BirationalDVR.logDiscrepancy_ne_top}
\leanok
\uses{lem:c0-dvr-primitive}
\source{fhx-envelope-20261028:lem:divisorial-trace-estimate}
In the setting of Lemma~\ref{lem:c0-dvr-primitive}, write $w_1=\pi$, $w_j=u_j$ for $2\le j\le N$. The partial derivatives
$\partial/\partial w_j$ of $L(\pi)=F(w_1,\dots,w_N)$ extend uniquely to $F$-derivations $\partial_j$ of $K$, and $\partial_j(V)\subseteq V$.
The Jacobian $\mathrm{Jac}_V\coloneqq\det(\partial_jz_i)_{i,j}\in V$ is nonzero. We put
$A(V)\coloneqq1+\operatorname{ord}_V(\mathrm{Jac}_V)\in\mathbf Z_{\ge1}$ (the log discrepancy of $\operatorname{ord}_V$ with respect to $z$, for the coordinates chosen).
\end{lemma}

\begin{proof}
For a derivation $D$ of $L(\pi)$ and $P\in L(\pi)[T]$ let $P^D$ be obtained by applying $D$ to the coefficients. Since $F_\theta$ is
separable, $F_\theta'(\theta)\ne0$, and $D$ extends uniquely to $K=L(\pi)[T]/(F_\theta)$ by $D(P(\theta))=P^D(\theta)+P'(\theta)D\theta$ with
$D\theta=-F_\theta^D(\theta)/F_\theta'(\theta)$; this is well defined because, for this choice of $D\theta$, the expression
$P^D(\theta)+P'(\theta)D\theta$ vanishes for $P=F_\theta Q$. The $\partial_j$ preserve $A$ (quotient rule; they act on $L=F(u)$ and on $\pi$), and $F_\theta'(\theta)$ is a unit of $V$ because its
residue $\bar F_\theta'(\xi)$ is nonzero (Lemma~\ref{lem:c0-dvr-primitive}(2)); so $\partial_j\theta\in V$, and by the quotient rule
$\partial_j(A[\theta]_{\mathfrak n'})\subseteq A[\theta]_{\mathfrak n'}=V$. Hence $\mathrm{Jac}_V\in V$. If $\mathrm{Jac}_V=0$, there is $0\ne(c_j)\in K^N$ with
$\sum_jc_j\partial_jz_i=0$ for all $i$; then $D=\sum_jc_j\partial_j$ is an $F$-derivation of $K=F(z)$ vanishing on the generators $z_i$, so
$D=0$, while $D(w_j)=c_j$, a contradiction.
\end{proof}

\subsection{The valuative bound}

In this subsection $F=\mathbf C$.

\begin{lemma}[Analytic coordinates at a general point of the divisor]\label{lem:c0-analytic-chart}
\lean{EnvelopeProperty.Char0.exists_analyticChart, EnvelopeProperty.Char0.exists_chartTheta, EnvelopeProperty.Char0.chartEv, EnvelopeProperty.Char0.eventually_det_chartMap, EnvelopeProperty.Char0.det_restrictScalars_real}
\leanok
\uses{lem:c0-dvr-primitive, lem:c0-dvr-jacobian}
Let $V$ be a divisorial valuation ring over $\mathbf C$ in $N$ variables and $f_1,\dots,f_p\in\mathbf C[z]$ nonzero. There are
$u_0\in\mathbf C^{N-1}$, $\delta\in(0,1)$ and a holomorphic map $\mu\colon U\to\mathbf C^N$ on
$U=\{(s,u)\in\mathbf C\times\mathbf C^{N-1}:|s|<\delta,\ |u-u_0|<\delta\}$, with bounded image, such that
\begin{enumerate}
\item $f_i\circ\mu=s^{\operatorname{ord}_V(f_i)}h_i$ with $h_i$ holomorphic and without zeros on $U$, for $1\le i\le p$;
\item $\det D\mu=s^{A(V)-1}h_0$ with $h_0$ holomorphic and without zeros on $U$;
\item $\mu$ is injective on $U\cap\{s\ne0\}$.
\end{enumerate}
\end{lemma}

\begin{proof}
Use the coordinates $\pi,u,\theta$ of Lemmas~\ref{lem:c0-dvr-curve}--\ref{lem:c0-dvr-jacobian}. Every $r\in V\setminus\{0\}$ is
$\pi^{\operatorname{ord}_V r}\varepsilon_r$ with $\varepsilon_r=P_r(\theta)/S_r(\theta)$, $P_r,S_r\in A[T]$ and $\bar P_r(\xi),\bar S_r(\xi)\ne0$ in $\kappa_V=L[T]/(G)$
(Lemma~\ref{lem:c0-dvr-primitive}(3)). Write $\pi=p_0(z)/p_1(z)$ and $u_j=q_j(z)/r_j(z)$ with polynomials, $p_1,r_j\ne0$, and
apply the following to the finite family $\mathcal F$ consisting of $f_1,\dots,f_p,p_1,r_2,\dots,r_N$ and $\mathrm{Jac}_V$.

\emph{Choice of $u_0$.} The finitely many elements of $L=\mathbf C(u)$ involved are: the coefficients (in $\pi$ and $T$) of the
numerators and denominators of the coefficients of $F_\theta,P_r,S_r$ ($r\in\mathcal F$) and of $z_i=P_{z_i}(\theta)/S_{z_i}(\theta)$; the
constant terms $S(0)\in L^\times$ of these denominators $S(\pi)$; the discriminant of $G$; and, for each $Q\in\{\bar P_r,\bar S_r,
\bar F_\theta/G\}$, the resultant $\operatorname{Res}_T(G,Q)\ne0$ ($G$ irreducible does not divide $Q$). Choose $u_0$ at which all of them are
regular and the nonzero ones do not vanish (the complement of the zero set of a nonzero polynomial is nonempty), and a root
$\theta_0$ of $G(u_0,T)$. Then $\theta_0$ is a simple root of $F_\theta(0,u_0,T)=\bar F_\theta(u_0,T)$, and $Q(u_0,\theta_0)\ne0$ for the $Q$
above.

\emph{The chart.} The coefficients of $F_\theta$ are holomorphic near $(0,u_0)$ and $\partial_TF_\theta(0,u_0,\theta_0)\ne0$, so the holomorphic
implicit function theorem gives a holomorphic $\vartheta(s,u)$ near $(0,u_0)$ with $\vartheta(0,u_0)=\theta_0$ and
$F_\theta(s,u,\vartheta(s,u))=0$. Evaluation $\pi\mapsto s$, $u\mapsto u$, $\theta\mapsto\vartheta$ defines a ring homomorphism $\mathrm{ev}$ from the
subring of $V$ of elements $P(\theta)/S(\theta)$ whose denominators $S$ are products of the $S_r$ and of the denominators above, to
holomorphic functions on a neighbourhood of $(0,u_0)$ (it is well defined: if $P(\theta)/S(\theta)=P'(\theta)/S'(\theta)$, then after clearing denominators $c(w)(PS'-P'S)=\Phi\cdot M$ with $0\ne c\in\mathbf C[w]$ and $\Phi$ the minimal polynomial of $\theta$ with cleared denominators, so $PS'-P'S$ vanishes on the graph of $\vartheta$ off the Lebesgue-null set $\{c=0\}$ (Lemma~\ref{lem:c0-analytic-basics}(4)), hence everywhere near $(0,u_0)$ by continuity; $F_\theta(\vartheta)=0$ alone does not suffice, because the coefficients of the quotient need not be regular at $(0,u_0)$). Put
$\mu=(\mathrm{ev}(z_1),\dots,\mathrm{ev}(z_N))$. For $r\in\mathcal F$, $\mathrm{ev}(r)=s^{\operatorname{ord}_Vr}\mathrm{ev}(\varepsilon_r)$ and
$\mathrm{ev}(\varepsilon_r)(0,u_0)=\bar P_r(u_0,\theta_0)/\bar S_r(u_0,\theta_0)\ne0$; since $\mathrm{ev}(f)=f\circ\mu$ for polynomials $f$, shrinking
$U$ gives (1). Differentiating $F_\theta(s,u,\vartheta)=0$ gives $\partial\vartheta/\partial w_j=\mathrm{ev}(\partial_j\theta)$, and by the chain rule
$\partial\,\mathrm{ev}(r)/\partial w_j=\mathrm{ev}(\partial_jr)$; hence $\det D\mu=\mathrm{ev}(\mathrm{Jac}_V)$, which gives (2). For (3), the identities
$p_1\pi=p_0$ and $r_ju_j=q_j$ in $V$ give $p_1(\mu(s,u))\,s=p_0(\mu(s,u))$ and $r_j(\mu(s,u))\,u_j=q_j(\mu(s,u))$; on $\{s\ne0\}$ the factors
$p_1\circ\mu$, $r_j\circ\mu$ have no zeros by (1), so $(s,u)$ is a function of $\mu(s,u)$. Shrinking $U$ makes $\mu(U)$ bounded.
\end{proof}

\begin{proposition}[Valuative bound for square integrable polynomials]\label{prop:c0-valuative-bound}
\lean{EnvelopeProperty.Char0.valuativeBound_of_coordinates, EnvelopeProperty.Char0.valuativeBound_forall_dvr, EnvelopeProperty.Char0.neg_one_lt_of_setLIntegral_rpow_lt_top}
\leanok
\uses{lem:c0-analytic-chart, lem:c0-analytic-basics}
\source{demailly-analytic-methods, fhx-envelope-20261028:lem:divisorial-trace-estimate}
Let $V$ be a divisorial valuation ring over $\mathbf C$ in $N$ variables, $g=(g_1,\dots,g_s)$ polynomials not all zero, $c>0$ real,
and $f\in\mathbf C[z]$ square integrable against $|g|^{-2c}$. Then $\operatorname{ord}_V(f)>c\operatorname{ord}_V(g)-A(V)$.
\end{proposition}

\begin{proof}
We may assume $f\ne0$ and drop the $g_j$ that vanish. Apply Lemma~\ref{lem:c0-analytic-chart} to $f,g_1,\dots,g_s$, shrink $U$ so that
the functions $h$ are bounded above and below on $U$ in modulus, and put $a=\operatorname{ord}_V f$, $b=\operatorname{ord}_V g=\min_j\operatorname{ord}_V g_j$,
$\alpha=A(V)-1$. On $U$: $|f\circ\mu|^2\ge c_1|s|^{2a}$, $|g\circ\mu|^2\le C_2|s|^{2b}$ (as $|s|<1$ and $\operatorname{ord}_Vg_j\ge b$), and $|g\circ\mu|>0$ on
$U_*=U\cap\{s\ne0\}$; the real Jacobian of the holomorphic map $\mu$ is $|\det D\mu|^2\ge c_3|s|^{2\alpha}$. Since $\mu$ is injective and $C^1$
on $U_*$, the change of variables formula and the hypothesis on the compact set $\overline{\mu(U)}$ give
\[
\infty>\int_{\mu(U_*)}|f|^2|g|^{-2c}\,d\lambda=\int_{U_*}|f\circ\mu|^2|g\circ\mu|^{-2c}|\det D\mu|^2\,d\lambda
\ge c_1c_3C_2^{-c}\int_{U_*}|s|^{2(a-cb+\alpha)}\,d\lambda(s,u).
\]
By Fubini and polar coordinates in $s$ (Lemma~\ref{lem:c0-analytic-basics}(3)) the last integral is finite only if $a-cb+\alpha>-1$,
that is, $\operatorname{ord}_Vf>c\operatorname{ord}_Vg-A(V)$.
\end{proof}

\begin{corollary}[Analytic lower bounds give valuative upper bounds]\label{cor:c0-valuative-upper}
\lean{EnvelopeProperty.Char0.ValuativeUpperBound, EnvelopeProperty.Char0.ValuativeUpperBound.of_valuative_bound, EnvelopeProperty.Char0.valuativeUpperBound_holds}
\leanok
\uses{prop:c0-valuative-bound}
Let $P=(P_1,\dots,P_r)$ and $G=(G_1,\dots,G_s)$ be polynomials, $G$ not all zero, and $a,b\in\mathbf Z_{>0}$. Suppose that for every
compact $K\subseteq\mathbf C^N$ there is $c_K>0$ with $|P|^b\ge c_K|G|^a$ on $K$. Then $b\operatorname{ord}_V(P)\le a\operatorname{ord}_V(G)$ for every divisorial
valuation ring $V$ over $\mathbf C$ in $N$ variables.
\end{corollary}

\begin{proof}
$P$ is not identically zero, since otherwise $G$ would vanish everywhere. Choose $j_0$ with $\operatorname{ord}_VG_{j_0}=\operatorname{ord}_VG$ and let
$\ell\ge1$ be an integer and $f=G_{j_0}^{a\ell}$. On a compact $K$, where $P\ne0$, $|f|^2|P|^{-2b\ell}\le|G|^{2a\ell}|P|^{-2b\ell}\le c_K^{-2\ell}$, and
where $P=0$ also $G=0$ and $f=0$; so $f$ is square integrable against $|P|^{-2b\ell}$. Proposition~\ref{prop:c0-valuative-bound} gives
$a\ell\operatorname{ord}_VG>b\ell\operatorname{ord}_VP-A(V)$; divide by $\ell$ and let $\ell\to\infty$.
\end{proof}

\subsection{\texorpdfstring{$L^2$}{L2} extension with polynomial growth}

\begin{lemma}[Standard facts from analysis]\label{lem:c0-analytic-basics}
\lean{EnvelopeProperty.Char0.taylorCoeff, EnvelopeProperty.Char0.torusIntegral_cauchy, EnvelopeProperty.Char0.hasSum_taylorCoeff, EnvelopeProperty.Char0.setIntegral_Icc_eq_taylorCoeff, EnvelopeProperty.Char0.volume_mvPolynomial_zeroSet_eq_zero, EnvelopeProperty.Char0.exists_differentiableOn_ae_eq_of_weakDbarEq_zero, EnvelopeProperty.Char0.not_integrableOn_ball_inv_norm_sub_sq_pow, EnvelopeProperty.Char0.eq_zero_of_lintegral_pole_ne_top, EnvelopeProperty.Char0.integrableOn_ball_norm_rpow_iff}
\leanok
\begin{enumerate}
\item (Weyl's lemma for $\bar\partial$.) A function $u\in L^2_{\mathrm{loc}}(\Omega)$ on an open $\Omega\subseteq\mathbf C^N$ with $\bar\partial u=0$ in the sense of
distributions agrees almost everywhere with a holomorphic function on $\Omega$.
\item (Cauchy's formula.) An entire function $F$ on $\mathbf C^N$ is the sum of its Taylor series $\sum_\gamma c_\gamma z^\gamma$, locally
uniformly, and $c_\gamma\rho^\gamma=(2\pi)^{-N}\int_{[0,2\pi]^N}F(\rho_1e^{i\vartheta_1},\dots,\rho_Ne^{i\vartheta_N})e^{-i\langle\gamma,\vartheta\rangle}\,d\vartheta$ for
$\rho\in\R_{>0}^N$.
\item $\int_{B(y,r)}|z-y|^{-2N}d\lambda=\infty$, and for $\delta>0$, $\kappa\in\R$: $\int_{\{|s|<\delta\}}|s|^{2\kappa}d\lambda(s)<\infty$ in $\mathbf C$ if and only if
$\kappa>-1$.
\item The zero set of a nonzero polynomial is $\lambda$-null.
\end{enumerate}
\end{lemma}

\begin{proof}
(1) Convolution with a mollifier commutes with $\bar\partial$; the regularizations $u_\epsilon$ are holomorphic and converge to $u$ in
$L^2_{\mathrm{loc}}$, and $L^2_{\mathrm{loc}}$ limits of holomorphic functions are holomorphic (mean value inequality). (2) Iterate the
one-variable Cauchy formula on polydiscs. (3) Polar coordinates: $\int_0^r\rho^{-2N}\rho^{2N-1}d\rho=\infty$ and
$\int_0^\delta\rho^{2\kappa+1}d\rho<\infty$ if and only if $\kappa>-1$. (4) Induction on $N$ with Fubini: for all but finitely many values of the
last variable the polynomial in the remaining variables is nonzero.
\end{proof}

\begin{lemma}[Weights of bihomogeneous families]\label{lem:c0-weight}
\lean{EnvelopeProperty.Char0.leviFormC0, EnvelopeProperty.Char0.leviFormC0_log_add_sum_normSq, EnvelopeProperty.Char0.le_leviFormC0_log_add_sum_normSq, EnvelopeProperty.Char0.normSqFam_xtRot, EnvelopeProperty.Char0.exists_normSqFam_le, EnvelopeProperty.Char0.rpow_neg_le_normSqFam_rpow_neg, EnvelopeProperty.Char0.volume_normSqFam_eq_zero, EnvelopeProperty.Char0.exists_normSqFam_le_on_compact, EnvelopeProperty.Char0.le_leviFormC0_log_one_add_norm_sq, EnvelopeProperty.Char0.le_leviFormC0_weight}
\leanok
\uses{conv:c0, lem:c0-analytic-basics}
Let $g=(g_1,\dots,g_s)$ be bihomogeneous polynomials of $x$-degree $D$ in $\mathbf C[x,t]$, not all zero, $c>0$, and
$\varphi=c\log|g|$.
\begin{enumerate}
\item For holomorphic $h_1,\dots,h_p$ and $\epsilon>0$, $\log(\sum_j|h_j|^2+\epsilon)$ is smooth and plurisubharmonic and decreases to
$\log\sum_j|h_j|^2$ as $\epsilon\downarrow0$; in particular $\varphi$ and $\log|z-y|$ are plurisubharmonic.
\item $|g(e^{i\vartheta}x,e^{i\eta}t)|=|g(x,t)|$ for real $\vartheta,\eta$.
\item There are $C_g>0$ and $\beta\in\N$ with $|g(x,t)|^2\le C_g(1+|x|^2)^D(1+|t|^2)^\beta$; hence
$e^{-2\varphi}\ge C_g^{-c}(1+|x|^2)^{-cD}(1+|t|^2)^{-c\beta}$.
\item $Z(g)$ is $\lambda$-null, and $\varphi$ is bounded above on compact sets.
\item $\log(1+|z|^2)$ is smooth with Levi form $\ge(1+|z|^2)^{-2}|\xi|^2$.
\end{enumerate}
\end{lemma}

\begin{proof}
(1) With $S=\sum|h_j|^2+\epsilon$ and $\partial h_j(\xi)=\sum_i\xi_i\partial h_j/\partial z_i$, the Levi form of $\log S$ is
$S^{-2}\bigl(S\sum_j|\partial h_j(\xi)|^2-|\sum_j\bar h_j\partial h_j(\xi)|^2\bigr)\ge0$ by the Cauchy--Schwarz inequality; decreasing limits of
plurisubharmonic functions are plurisubharmonic. (2) $g_j(e^{i\vartheta}x,e^{i\eta}t)=e^{iD\vartheta+i\lambda_j\eta}g_j(x,t)$ for $g_j=s_j(x)t^{\lambda_j}$. (3)
Each monomial $x^\alpha t^\lambda$ of $g_j$ has $|\alpha|=D$, so $|x^\alpha t^\lambda|\le(1+|x|^2)^{D/2}(1+|t|^2)^{\beta/2}$ with $\beta$ the largest
$t$-degree. (4) $Z(g)\subseteq Z(g_j)$ for $g_j\ne0$ (Lemma~\ref{lem:c0-analytic-basics}(4)), and $|g|$ is continuous. (5) The Levi form is
$(1+|z|^2)^{-2}\bigl((1+|z|^2)|\xi|^2-|\langle\xi,z\rangle|^2\bigr)$.
\end{proof}

\begin{lemma}[The Ohsawa--Takegoshi theorem at a point]\label{lem:c0-ot-point}
\lean{EnvelopeProperty.Char0.exists_extension_from_point_ball, EnvelopeProperty.Char0.plurisubharmonicOn_mul_logNormSqFam}
\leanok
\source{demailly-analytic-methods, demailly92-regularization}
For every $r>0$ there is $C_r>0$ such that for every $y\in\mathbf C^N$ and every plurisubharmonic $\psi$ on $B(y,r)$ with $\psi(y)>-\infty$
there is a holomorphic function $f$ on $B(y,r)$ with $f(y)=1$ and $\int_{B(y,r)}|f|^2e^{-\psi}\,d\lambda\le C_re^{-\psi(y)}$.
\end{lemma}

\begin{proof}
This is the Ohsawa--Takegoshi extension theorem from the point $\{y\}$ in the bounded pseudoconvex domain $B(y,r)$,
\cite[Corollary~13.9, Remark~13.10]{DemAM}, used in the same way in the proof of \cite[Theorem~14.2]{DemAM}; translate to $y=0$.
\end{proof}

\begin{lemma}[H\"ormander's estimate on balls]\label{lem:c0-hormander-bounded}
\lean{EnvelopeProperty.Char0.hormander_ball, EnvelopeProperty.Char0.exists_holomorphic_ball}
\leanok
\source{demailly-analytic-methods}
Let $R>0$, $\Phi$ a smooth plurisubharmonic function on a neighbourhood of $\overline{B(0,R+1)}\subseteq\mathbf C^N$, $v=\sum_jv_jd\bar z_j$ a
smooth $(0,1)$-form with compact support in $B(0,R)$ and $\bar\partial v=0$, and $g\ge0$ measurable, supported in $B(0,R)$, with
$|\sum_jv_j\bar\xi_j|^2\le g\cdot L_\Phi(\xi)$ for all $\xi\in\mathbf C^N$ ($L_\Phi$ the Levi form) and $\int g e^{-\Phi}<\infty$. There is
$u\in L^2_{\mathrm{loc}}(B(0,R))$ with $\bar\partial u=v$ on $B(0,R)$ and $\int_{B(0,R)}|u|^2e^{-\Phi}\,d\lambda\le2\int ge^{-\Phi}\,d\lambda$.
\end{lemma}

\begin{proof}
This is H\"ormander's $L^2$ estimate for the $\bar\partial$-equation on the bounded convex domain $B(0,R+1)$ with the weight
$e^{-\Phi}$, in the twisted form with the Levi form of the weight \cite[Theorems~5.1 and 5.2]{DemAM}, restricted to $B(0,R)$.
\end{proof}

\begin{lemma}[H\"ormander's estimate on $\mathbf C^N$ with polynomial weights]\label{lem:c0-hormander-global}
\lean{EnvelopeProperty.Char0.exists_entire_hormander_global, EnvelopeProperty.Char0.exists_entire_hormander_global_weights, EnvelopeProperty.Char0.hormanderWeight, EnvelopeProperty.Char0.iSup_exp_neg_hormanderWeight}
\leanok
\uses{lem:c0-hormander-bounded, lem:c0-weight}
\source{demailly-analytic-methods}
Let
\[
 \Phi_0=\sum_{a=1}^pc_a\log\sum_j|h_{a,j}|^2,\qquad c_a>0,
\]
where the $h_{a,j}$ are entire and not all identically zero for each $a$. Set
$\Phi=\Phi_0+\log(1+|z|^2)$. Let $v$ be a smooth $(0,1)$-form with compact support, $\bar\partial v=0$, and
$I=\int|v|^2(1+|z|^2)^2e^{-\Phi}d\lambda<\infty$. Then there is $u\in L^2_{\mathrm{loc}}(\mathbf C^N)$ with $\bar\partial u=v$ and
$\int_{\mathbf C^N}|u|^2e^{-\Phi}d\lambda\le2I$.
\end{lemma}

\begin{proof}
For $\epsilon\in(0,1]$ let $\Phi_\epsilon=\sum_ac_a\log(\sum_j|h_{a,j}|^2+\epsilon)+\log(1+|z|^2)$, smooth plurisubharmonic, decreasing to $\Phi$
(Lemma~\ref{lem:c0-weight}(1)). By Lemma~\ref{lem:c0-weight}(1),(5), $L_{\Phi_\epsilon}(\xi)\ge(1+|z|^2)^{-2}|\xi|^2$, so
$|\sum v_j\bar\xi_j|^2\le g\,L_{\Phi_\epsilon}(\xi)$ with $g=|v|^2(1+|z|^2)^2$, independent of $\epsilon$, and $\int ge^{-\Phi_\epsilon}\le I$ since
$\Phi_\epsilon\ge\Phi$. Let $\operatorname{supp}v\subseteq B(0,R_0)$. For integers $R\ge R_0$, Lemma~\ref{lem:c0-hormander-bounded} gives $u_{R,\epsilon}$ on $B(0,R)$
with $\bar\partial u_{R,\epsilon}=v$ and $\int_{B(0,R)}|u_{R,\epsilon}|^2e^{-\Phi_\epsilon}\le2I$. On each ball $B(0,R')$ the weights $e^{-\Phi_\epsilon}$ are bounded
above and below, so for fixed $\epsilon$ the restrictions are bounded in $L^2(B(0,R'))$; a diagonal subsequence converges weakly on every
ball to some $u_\epsilon$, with $\bar\partial u_\epsilon=v$ (test against compactly supported forms) and, by weak lower semicontinuity of the norm of
$L^2(B(0,R'),e^{-\Phi_\epsilon}\lambda)$, $\int_{B(0,R')}|u_\epsilon|^2e^{-\Phi_\epsilon}\le2I$ for all $R'$, hence on $\mathbf C^N$. For $\epsilon\le\epsilon'$,
$e^{-\Phi_{\epsilon'}}\le e^{-\Phi_\epsilon}$, so $\int|u_\epsilon|^2e^{-\Phi_{\epsilon'}}\le2I$; again a diagonal subsequence $\epsilon=1/\ell\to0$ converges weakly on every ball to
some $u$ with $\bar\partial u=v$ and $\int|u|^2e^{-\Phi_{\epsilon'}}\le2I$ for every $\epsilon'$. Monotone convergence as $\epsilon'\downarrow0$ gives
$\int|u|^2e^{-\Phi}\le2I$.
\end{proof}

\begin{lemma}[Square integrable entire functions of polynomial growth]\label{lem:c0-polynomial-growth}
\lean{EnvelopeProperty.Char0.exists_mvPolynomial_of_lintegral_polyWeight, EnvelopeProperty.Char0.polyWeight_add_le}
\leanok
\uses{lem:c0-analytic-basics}
Let $F$ be an entire function on $\mathbf C^{n+1}_x\times\mathbf C_t$ and $a,b\ge0$ with
$\|F\|^2\coloneqq\int|F|^2(1+|x|^2)^{-a}(1+|t|^2)^{-b}d\lambda<\infty$. Then $F$ is a polynomial with $x$-degree $\le a-n-1$ and $t$-degree
$\le b-1$.
\end{lemma}

\begin{proof}
Write $F=\sum c_{\alpha,\lambda}x^\alpha t^\lambda$ (Lemma~\ref{lem:c0-analytic-basics}(2)). For $\rho,\sigma\ge1$ let
$S=\{\rho\le|x_i|\le2\rho\ \forall i,\ \sigma\le|t|\le2\sigma\}$. Integrating Cauchy's formula over the radii in polar coordinates,
\[
|c_{\alpha,\lambda}|\rho^{|\alpha|}\sigma^\lambda\Bigl(\tfrac32\rho^2\Bigr)^{n+1}\tfrac32\sigma^2(2\pi)^N\le\int_S|F|\,d\lambda\le
\lambda(S)^{1/2}\Bigl(\sup_S(1+|x|^2)^a(1+|t|^2)^b\Bigr)^{1/2}\|F\|,
\]
Here
\[
 \lambda(S)\le C\rho^{2(n+1)}\sigma^2,\qquad
 \sup_S(1+|x|^2)^a(1+|t|^2)^b\le C\rho^{2a}\sigma^{2b}.
\]
Hence
$|c_{\alpha,\lambda}|\le C'\rho^{a-n-1-|\alpha|}\sigma^{b-1-\lambda}\|F\|$ for all $\rho,\sigma\ge1$. Letting $\rho\to\infty$ (resp.\ $\sigma\to\infty$) shows
$c_{\alpha,\lambda}=0$ if $|\alpha|>a-n-1$ (resp.\ $\lambda>b-1$).
\end{proof}

\begin{proposition}[$L^2$ extension from a point with polynomial growth]\label{prop:c0-l2-extension}
\lean{EnvelopeProperty.Char0.exists_entire_extension, EnvelopeProperty.Char0.exists_polynomial_extension, EnvelopeProperty.Char0.c0Weight}
\leanok
\uses{lem:c0-ot-point, lem:c0-hormander-global, lem:c0-polynomial-growth, lem:c0-weight, lem:c0-analytic-basics}
\source{demailly-analytic-methods, demailly92-regularization}
Let $g$, $D$, $c$, $\varphi$, $C_g$, $\beta$ be as in Lemma~\ref{lem:c0-weight}, $M=cD$, and $\omega=e^{-2\varphi}(1+|z|^2)^{-N-1}$. For every compact
$K\subseteq\mathbf C^N$ there is $C_K$ such that for every $y\in K\setminus Z(g)$ there is a polynomial $F$ with $F(y)=1$,
$\int|F|^2\omega\,d\lambda\le C_Ke^{-2\varphi(y)}$, $x$-degree $\le M+2$ and $t$-degree $\le c\beta+N$.
\end{proposition}

\begin{proof}
Let $R_K=\max_K|z|$ and $y\in K\setminus Z(g)$. (1) Lemma~\ref{lem:c0-ot-point} with $r=1$ and $\psi=2\varphi$ gives $f$ holomorphic on $B(y,1)$ with
$f(y)=1$ and $\int_{B(y,1)}|f|^2e^{-2\varphi}\le C_1e^{-2\varphi(y)}$. (2) Fix a smooth $\chi\colon\mathbf C^N\to[0,1]$ equal to $1$ on $B(0,1/2)$ with
support in $B(0,1)$, let $\chi_y=\chi(\cdot-y)$ and $v=f\,\bar\partial\chi_y$, extended by $0$: a smooth $\bar\partial$-closed $(0,1)$-form supported in
$\{1/2\le|z-y|\le1\}$. (3) Let $\Phi_0=2\varphi+2N\log|z-y|=c\log\sum_j|g_j|^2+N\log\sum_i|z_i-y_i|^2$ and $\Phi=\Phi_0+\log(1+|z|^2)$. On the
support of $v$, $e^{-\Phi}\le2^{2N}e^{-2\varphi}$, so $I=\int|v|^2(1+|z|^2)^2e^{-\Phi}\le C_2e^{-2\varphi(y)}$ with
$C_2=2^{2N}\sup|\bar\partial\chi|^2(1+(R_K+1)^2)^2C_1$. Lemma~\ref{lem:c0-hormander-global} gives $u$ with $\bar\partial u=v$ and
$\int|u|^2e^{-\Phi}\le2C_2e^{-2\varphi(y)}$. (4) $F\coloneqq\chi_yf-u$ has $\bar\partial F=0$, so it is (almost everywhere equal to) an entire
function (Lemma~\ref{lem:c0-analytic-basics}(1)). On $B(y,1/2)$, $u=f-F$ is holomorphic, and $\varphi$ is bounded above there
(Lemma~\ref{lem:c0-weight}(4)), so $\int_{B(y,1/2)}|u|^2|z-y|^{-2N}<\infty$; by Lemma~\ref{lem:c0-analytic-basics}(3) and continuity, $u(y)=0$. Hence
$F(y)=1$. (5) Since $|z-y|^2\le2(1+R_K^2)(1+|z|^2)$, $\omega\le(2(1+R_K^2))^Ne^{-\Phi}$, and $\int|\chi_yf|^2\omega\le C_1e^{-2\varphi(y)}$; so
$\int|F|^2\omega\le C_Ke^{-2\varphi(y)}$ with $C_K=2C_1+4(2(1+R_K^2))^NC_2$. (6) By Lemma~\ref{lem:c0-weight}(3) and
$1+|z|^2\le(1+|x|^2)(1+|t|^2)$, $\omega\ge C_g^{-c}(1+|x|^2)^{-M-N-1}(1+|t|^2)^{-c\beta-N-1}$, and Lemma~\ref{lem:c0-polynomial-growth} with $a=M+N+1$,
$b=c\beta+N+1$ shows that $F$ is a polynomial of $x$-degree $\le M+N+1-(n+1)=M+2$ and $t$-degree $\le c\beta+N$.
\end{proof}

\begin{corollary}[Polynomials dominating a fractional power]\label{cor:c0-bergman-lower}
\lean{EnvelopeProperty.Char0.bergmanSpace, EnvelopeProperty.Char0.wNormSq_weightedHomogeneousComponent_le, EnvelopeProperty.Char0.exists_homogeneous_basis, EnvelopeProperty.Char0.exists_sum_norm_le, EnvelopeProperty.Char0.exists_homogeneous_lower_bound, EnvelopeProperty.Char0.exists_homogeneous_lower_bound_of_bihomogeneous, EnvelopeProperty.Char0.BergmanLowerBound, EnvelopeProperty.Char0.bergmanLowerBound_holds}
\leanok
\uses{prop:c0-l2-extension, lem:c0-weight, lem:c0-analytic-basics}
\source{demailly-analytic-methods}
In the setting of Proposition~\ref{prop:c0-l2-extension} there are finitely many $x$-homogeneous polynomials $P_1,\dots,P_r$ of
$x$-degree $\le M+2$, each square integrable against $|g|^{-2c}$, such that for every compact $K$ there is $c_K>0$ with
$|P|\ge c_K|g|^c$ on $K$.
\end{corollary}

\begin{proof}
Let $W$ be the space of polynomials $F$ of $x$-degree $\le M+2$ and $t$-degree $\le c\beta+N$ with $\|F\|_\omega^2=\int|F|^2\omega\,d\lambda<\infty$. It is a
finite dimensional vector space, and $\|\cdot\|_\omega$ is a norm on it ($\omega>0$ almost everywhere). The measure $\omega\lambda$ is invariant under
$x\mapsto e^{i\vartheta}x$ (Lemma~\ref{lem:c0-weight}(2)), so for $F\in W$ the $x$-homogeneous components
$F_e=(2\pi)^{-1}\int_0^{2\pi}e^{-ie\vartheta}F(e^{i\vartheta}x,t)\,d\vartheta$ satisfy $\|F_e\|_\omega\le\|F\|_\omega$ (Minkowski's inequality) and lie in $W$. Let
$P_1,\dots,P_r$ be a basis of $W$ consisting of $x$-homogeneous polynomials. Each $P_k$ is square integrable against $|g|^{-2c}=e^{-2\varphi}$,
since $\int_K|P_k|^2e^{-2\varphi}\le\sup_K(1+|z|^2)^{N+1}\|P_k\|_\omega^2$. By equivalence of norms on $W$ there is $C_W$ with
$\sum|a_k|\le C_W\|\sum a_kP_k\|_\omega$. For $y\in K\setminus Z(g)$ let $F_y=\sum_ka_kP_k\in W$ be given by Proposition~\ref{prop:c0-l2-extension}; then
$1=|F_y(y)|\le\sum_k|a_k|\,|P(y)|\le C_WC_K^{1/2}e^{-\varphi(y)}|P(y)|$, so $|P(y)|\ge c_K|g(y)|^c$ with $c_K=(C_WC_K^{1/2})^{-1}$. On $Z(g)$ the
inequality is trivial.
\end{proof}

\subsection{Assembly}

In this subsection $k$, $\iota\colon k\to\mathbf C$ (Lemma~\ref{lem:c0-embedding}), $d\in\mathbf Z_{>0}$ and a rational PL function $\phi\le0$ on
$(\PP^n)^\an$ are fixed; $\phi=\psi_1-\psi_2$ with $\psi_\nu\in\mathcal H_\Q(b_\nu)$, and $J_m=J_m(\phi)$ (Definition~\ref{def:bj-Jm}). For
$h\in k[z]$ let $h^\iota\in\mathbf C[z]$ be its image, and for a divisorial point $v$ with ring $V$ (Definition~\ref{def:c0-dvr}) let $V_{\mathbf C}$
be its extension of scalars along $\iota$ (Lemma~\ref{lem:c0-base-change-dvr}).

\begin{lemma}[Finitely many valuations define $J_m$]\label{lem:c0-shilov-description}
\lean{EnvelopeProperty.ProjBerk.IsRatPL.exists_shilov_description}
\leanok
\uses{lem:bj-shilov, def:bj-Jm, lem:bj-fs-basic, def:c0-dvr}
\source{bj22-trivval:lem:Shilov}
There are divisorial points $v_1,\dots,v_r$, with rings $V_i$ and constants $c_i\in\Q_{>0}$ as in Definition~\ref{def:c0-dvr}, and real numbers
$\gamma_1,\dots,\gamma_r$ such that for every $m\ge1$,
\[
V_m(\phi)=\{g\in R_{md}[t]:\operatorname{ord}_{V_i}(g)\ge m\gamma_i\text{ for }1\le i\le r\}.
\]
\end{lemma}

\begin{proof}
Let $\Sigma=\{v_1,\dots,v_r\}$ be given by Lemma~\ref{lem:bj-shilov} for $\psi_1$, and $\gamma_i=-\phi(v_i)/c_i$ ($\phi$ is finite valued). The element
$g=0$ lies in both sets. For $0\ne g=\sum_\lambda s_\lambda t^\lambda\in R_{md}[t]$, $m^{-1}\log|g|_{\sigma(\cdot)}=m^{-1}\max_\lambda(\log|s_\lambda|-\lambda)$ is a
Fubini--Study function for $\mathcal O(d)$, so $\varphi_g=m^{-1}\log|g|_{\sigma(\cdot)}+\psi_2\in\mathcal H^{\mathrm{gf}}_\Q(d+b_2)$
(Lemma~\ref{lem:bj-fs-basic}). Now $g\in V_m(\phi)$ if and only if $\varphi_g\le\psi_1$ everywhere, if and only if (Lemma~\ref{lem:bj-shilov})
$\varphi_g\le\psi_1$ on $\Sigma$, if and only if $\log|g|_{\sigma(v_i)}\le m\phi(v_i)$ for all $i$, that is, $-c_i\operatorname{ord}_{V_i}(g)\le m\phi(v_i)$.
\end{proof}

\begin{proposition}[The key inequality]\label{prop:c0-key-inequality}
\lean{EnvelopeProperty.ProjBerk.IsRatPL.exists_key_inequality}
\leanok
\uses{lem:c0-shilov-description, cor:c0-bergman-lower, prop:c0-valuative-bound, cor:c0-valuative-upper, lem:c0-base-change-dvr, cor:c0-base-change-orders, lem:c0-dvr-jacobian, lem:c0-embedding, lem:bj-Jm-basic, thm:bj-divisorial-comparison, thm:bj-psh-basic}
\source{bj22-trivval:thm:contenvlisse, fhx-envelope-20261028:eq:smooth-Q-test-sandwich}
There is $N_0\in\N$ such that for all $m,q\ge1$, $\frac mq\phi_{J_q}\le\phi_{J_{m+2}}+N_0$ on $(\PP^n)^\an$.
\end{proposition}

\begin{proof}
Let $v_i,V_i,c_i,\gamma_i$ be as in Lemma~\ref{lem:c0-shilov-description}, $A_i=A(V_{i,\mathbf C})$ (Lemma~\ref{lem:c0-dvr-jacobian} for the divisorial
valuation ring $V_{i,\mathbf C}$ over $\mathbf C$), and $\tau_i=\operatorname{ord}_{V_i}(t)=1/c_i>0$ (as $\log|t|_{\sigma(v_i)}=-1$). Put
$N_0=\max_i\lceil(A_i+2|\gamma_i|)/\tau_i\rceil$.

Fix $m,q\ge1$. By Lemma~\ref{lem:bj-Jm-basic}, $J_q$ is generated by finitely many bihomogeneous $g_1,\dots,g_s$ of $x$-degree $qd$, not all zero.
Apply Corollary~\ref{cor:c0-bergman-lower} to $g^\iota=(g_1^\iota,\dots,g_s^\iota)$, $D=qd$ and $c=m/q$, so $M=md$: we get $x$-homogeneous
$P_1,\dots,P_r\in\mathbf C[z]$ of $x$-degree $\le md+2$, square integrable against $|g^\iota|^{-2m/q}$, with $|P|^q\ge c_K^q|g^\iota|^m$ on compact
sets $K$.

\emph{Step 1: $t^{N_0}x^\alpha P_k\in J_{m+2}\mathbf C[z]$ whenever $|\alpha|=(m+2)d-\deg_xP_k$.} By Proposition~\ref{prop:c0-valuative-bound} at $V_{i,\mathbf C}$ and
Corollary~\ref{cor:c0-base-change-orders}(2), $\operatorname{ord}_{V_{i,\mathbf C}}(P_k)>\frac mq\operatorname{ord}_{V_i}(J_q)-A_i\ge m\gamma_i-A_i$, because every generator
$g_j\in V_q(\phi)$ has $\operatorname{ord}_{V_i}(g_j)\ge q\gamma_i$ (Lemma~\ref{lem:c0-shilov-description}). As $x_l,t\in V_i$,
$G=t^{N_0}x^\alpha P_k$ has $\operatorname{ord}_{V_{i,\mathbf C}}(G)\ge N_0\tau_i+m\gamma_i-A_i\ge(m+2)\gamma_i$ for every $i$. $G$ is $x$-homogeneous of degree
$(m+2)d$; let $W\subseteq R_{(m+2)d}[t]$ be the finite dimensional subspace of $t$-degree $\le\deg_tG$. By Lemma~\ref{lem:c0-shilov-description} and
Corollary~\ref{cor:c0-base-change-orders}(1), $G\in(W\cap V_{m+2}(\phi))\otimes_k\mathbf C\subseteq J_{m+2}\mathbf C[z]$.

\emph{Step 2: divisorial points.} Let $v\in(\PP^n)^{\divv}$ with ring $V$ and constant $c_v$. As $\max_l|x_l|_v=1$, some $x_l$ has
$\operatorname{ord}_V(x_l)=0$, hence $\operatorname{ord}_{V_{\mathbf C}}(x_l^e)=0$. By Step~1 and Corollary~\ref{cor:c0-base-change-orders}(2), for each $k$,
\[
\operatorname{ord}_V(J_{m+2})=\operatorname{ord}_{V_{\mathbf C}}(J_{m+2}\mathbf C[z])\le N_0\operatorname{ord}_V(t)+\operatorname{ord}_{V_{\mathbf C}}(P_k).
\]
Taking the minimum over $k$ and applying Corollary~\ref{cor:c0-valuative-upper} with $(a,b)=(m,q)$ at $V_{\mathbf C}$:
$\operatorname{ord}_V(J_{m+2})\le N_0\operatorname{ord}_V(t)+\frac mq\operatorname{ord}_{V_{\mathbf C}}(g^\iota)=N_0\operatorname{ord}_V(t)+\frac mq\operatorname{ord}_V(J_q)$. Multiplying by $-c_v$, and using
$c_v\operatorname{ord}_V(t)=1$ and $\phi_I(v)=-c_v\operatorname{ord}_V(I)$, we get $\phi_{J_{m+2}}(v)+N_0\ge\frac mq\phi_{J_q}(v)$.

\emph{Step 3.} $u=\frac mq\phi_{J_q}$ is $md$-psh ($q^{-1}\phi_{J_q}\in\mathcal H_\Q(d)$ by Lemma~\ref{lem:bj-Jm-basic}, and Theorem~\ref{thm:bj-psh-basic}), and
$h=\phi_{J_{m+2}}+N_0$ is continuous and real valued. Since $u\le h$ on $(\PP^n)^{\divv}$ (Step~2), Theorem~\ref{thm:bj-divisorial-comparison} gives
$u\le h$ everywhere.
\end{proof}

\begin{theorem}[Envelopes of PL obstacles in characteristic zero]\label{thm:c0-envelope-pl}
\lean{EnvelopeProperty.ProjBerk.IsRatPL.isPSH_envelope_nat_of_nonpos_of_inputs}
\leanok
\uses{prop:c0-key-inequality, prop:bj-envelope-Jm, lem:bj-Jm-basic, thm:bj-psh-basic}
\source{bj22-trivval:lem:envline, bj22-trivval:thm:contenvlisse, fhx-envelope-20261028:eq:smooth-Q-limit}
In the setting of this subsection, $m^{-1}\phi_{J_m}$ converges uniformly to $\mathrm P_d(\phi)$ on $(\PP^n)^\an$, and $\mathrm P_d(\phi)\in\PSH(d)$.
\end{theorem}

\begin{proof}
$\phi$ is continuous, bounded and $\le0$, so $\mathrm Q\coloneqq\mathrm P_d(\phi)=\sup_qq^{-1}\phi_{J_q}$ (Proposition~\ref{prop:bj-envelope-Jm}). Taking the
supremum over $q$ in Proposition~\ref{prop:c0-key-inequality}, $m\mathrm Q\le\phi_{J_{m+2}}+N_0$ for $m\ge1$. Since $\mathrm Q\le\phi\le0$, for $m'\ge3$
\[
\mathrm Q-\frac{N_0}{m'}\le\frac{(m'-2)\mathrm Q-N_0}{m'}\le\frac1{m'}\phi_{J_{m'}}\le\mathrm Q,
\]
so $m'^{-1}\phi_{J_{m'}}\to\mathrm Q$ uniformly. These functions lie in $\mathcal H_\Q(d)\subseteq\PSH(d)$ (Lemma~\ref{lem:bj-Jm-basic}), and
$\PSH(d)$ is closed under uniform limits of sequences (Theorem~\ref{thm:bj-psh-basic}(2)), so $\mathrm Q\in\PSH(d)$.
\end{proof}

\begin{theorem}[Envelope property of $\PP^n$ in characteristic zero]\label{thm:c0-pn-ep}
\lean{EnvelopeProperty.ProjBerk.hasEnvelopeProperty_of_charZero, EnvelopeProperty.ProjBerk.hasEnvelopeProperty_of_charZero_of_inputs, EnvelopeProperty.ProjBerk.hasEnvelopeProperty_of_charZero_of_valuativeUpperBound, EnvelopeProperty.ProjBerk.hasEnvelopeProperty_of_ringHom_complex_of_inputs}
\leanok
\uses{thm:c0-envelope-pl, lem:c0-embedding, lem:bj-envppty, prop:bj-envelope-criterion, lem:bj-env-basic, lem:bj-pl-lattice, thm:bj-psh-basic, def:pn-ep}
\source{fhx-envelope-20261028:thm:smooth-EP, bj22-trivval:thm:contenvlisse}
Let $k$ be a countable algebraically closed field of characteristic $0$. Every $\theta\ge0$ has the envelope property on $\PP^n$.
\end{theorem}

\begin{proof}
By Lemma~\ref{lem:bj-envppty} applied to rationals $\theta_j\downarrow\theta$ with $\theta_j>\theta$, it suffices to treat $\theta=a\in\Q_{>0}$. By
Proposition~\ref{prop:bj-envelope-criterion} it suffices that $\mathrm P_a(\phi)$ is $a$-psh for every $\phi\in\operatorname{PL}$. Choose $m\ge1$ with
$d=ma\in\mathbf Z$ and a rational constant $c\ge\max m\phi$; by Lemma~\ref{lem:bj-env-basic},
$\mathrm P_a(\phi)=m^{-1}(\mathrm P_d(m\phi-c)+c)$, and $m\phi-c\le0$ is PL (Lemma~\ref{lem:bj-pl-lattice}). Theorem~\ref{thm:c0-envelope-pl} (with the embedding of
Lemma~\ref{lem:c0-embedding}) shows that $\mathrm P_d(m\phi-c)$ is $d$-psh, hence $\mathrm P_a(\phi)$ is $a$-psh (Theorem~\ref{thm:bj-psh-basic}).
\end{proof}

\begin{remark}[Relation to the sources]\label{rmk:c0-source}
\uses{thm:c0-pn-ep, rmk:bj-char-zero, rmk:bj-pn-ep-assembly}
\cite[Theorem~5.20]{BJ22} in characteristic $0$ passes to a resolution $\mathcal X$ of a test configuration $(\mathcal X,\mathcal L)$ with
$\phi=\phi_{\mathcal L}$, and compares the base ideals $\mathfrak a_m$ of $m\mathcal L$ with the asymptotic multiplier ideals
$\mathfrak b_m=\mathcal J(\mathfrak a_\bullet^m)$: $\mathfrak a_m\subseteq\mathfrak b_m$, subadditivity $\mathfrak b_{m\ell}\subseteq\mathfrak b_m^\ell$, and uniform global
generation of $\mathfrak b_m\otimes\mathcal O(m\mathcal L+m_0\mathcal A)$. Here the ideals $J_m$ on the affine cone replace the test
configuration (Remark~\ref{rmk:bj-test-configurations}), as in the characteristic $p$ argument of \cite[\S3]{FHX}, and multiplier ideals
are not used as ideals at all. Only the finite dimensional spaces of square integrable polynomials of bounded degree of
Corollary~\ref{cor:c0-bergman-lower} enter. The two properties of $\mathfrak b_m$ used by Boucksom--Jonsson become:
\begin{itemize}
\item subadditivity, which gives $m^{-1}\phi_{\mathfrak b_m}\ge\mathrm P$, is replaced by the lower bound $|P|\ge c_K|J_q|^{m/q}$ from the
Ohsawa--Takegoshi theorem at a point (Demailly's estimate $\varphi_m\ge\varphi-C/m$ for Bergman kernels, \cite[Theorem~14.2]{DemAM}), turned
into $\operatorname{ord}_V(P)\le\frac mq\operatorname{ord}_V(J_q)$ by Corollary~\ref{cor:c0-valuative-upper};
\item uniform global generation (Nadel vanishing and Castelnuovo--Mumford regularity, with $K_{\PP^n}+(n+1)H=0$) is replaced by
H\"ormander's estimate on $\mathbf C^{n+2}$ with polynomial weights, which bounds the $x$-degree by $md+2$ instead of $md$; the loss
of $2$ is absorbed by passing from $J_m$ to $J_{m+2}$.
\end{itemize}
The valuative bound $\operatorname{ord}_V(\mathcal J(I^c))\ge c\operatorname{ord}_V(I)-A(V)$, the analogue of \cite[Lemma~3.3]{FHX}, is only needed in the
direction of Proposition~\ref{prop:c0-valuative-bound}, which is computed at a general point of the divisor of $V$ (the change of
variables formula of \cite[Proposition~5.8]{DemAM} in a single chart) and needs no resolution of singularities. The constant $N_0$ of
Proposition~\ref{prop:c0-key-inequality} comes from the log discrepancies of finitely many divisorial valuations, as the constant
$A_S(w)$ in \cite[\S3]{FHX}. Countability of $k$ is used only to embed $k$ into $\mathbf C$.
\end{remark}
